\section*{Bab \ref{ch:g6:life-through-seasons} --- Kehidupan Sepanjang Musim}

\begin{solution}{exo:g6:life-through-seasons:1}
Panjang hari (\omterm{def:g6:exploring-our-environment:conditions}{cahaya}), \omterm{def:g6:exploring-our-environment:conditions}{suhu}, dan \omterm{def:g6:exploring-our-environment:conditions}{kelembapan} tanah atau udara yang
mengikuti hujan serta kemarau.
\end{solution}

\begin{solution}{exo:g6:life-through-seasons:2}
\omterm{def:g6:life-through-seasons:trees}{Meranggas} --- menanggalkan seluruh \omterm{def:g1:plants-around-us:parts}{daunnya} pada musim gugur: pohon ek,
beech, berangan kuda. \omterm{def:g6:life-through-seasons:trees}{Selalu hijau} --- mempertahankan \omterm{def:g1:plants-around-us:parts}{daun} yang
bekerja sepanjang musim dingin: pinus, holi.
\end{solution}

\begin{solution}{exo:g6:life-through-seasons:3}
Sebuah tunas lengkap yang mungil --- \omterm{def:g1:plants-around-us:parts}{daun} terlipat yang mungil,
terbungkus \omterm{def:g1:animals-around-us:covering}{sisik} yang tahan cuaca. Ia dibangun pada musim panas
sebelumnya.
\end{solution}

\begin{solution}{exo:g6:life-through-seasons:4}
\omterm{prop:g6:life-through-seasons:herbs}{Tumbuhan semusim} mati seluruhnya lalu melewati musim dingin sebagai
\omterm{def:g2:from-seed-to-plant:seed}{biji}. Bagian atas \omterm{def:g3:flowers-fruits-seeds:parts}{bunga} narsisnya mati sementara \omterm{def:g1:plants-around-us:plant}{tumbuhannya} menunggu
di bawah tanah sebagai \omterm{ex:g4:plants-without-seeds:bulbtuber}{umbi lapis}.
\end{solution}

\begin{solution}{exo:g6:life-through-seasons:5}
Karena tunas di dalam \omterm{def:g6:life-through-seasons:buds}{kuncupnya} sudah selesai dibangun pada musim
panas lalu; \omterm{def:g1:plants-around-us:tree}{rantingnya} hanya kekurangan kehangatan, yang disediakan
air di dalam ruangan berminggu-minggu sebelum kebunnya menyediakannya.
\end{solution}

\begin{solution}{exo:g6:life-through-seasons:6}
\omterm{def:g3:animals-through-seasons:migration}{Bermigrasi}: \omterm{prop:g3:sorting-living-things:groups}{burung} layang-layang. \omterm{def:g3:animals-through-seasons:hibernation}{Berhibernasi}: landak (atau siput
yang tersegel). \omterm{def:g3:animals-through-seasons:active}{Tetap aktif}: rubah, \omterm{prop:g3:sorting-living-things:groups}{burung} murai, \omterm{prop:g3:sorting-living-things:groups}{burung} robin.
Melewati musim dingin sebagai tahap muda --- \omterm{def:g2:animals-grow-up:born}{telur}, larva atau
\omterm{ex:g2:animals-grow-up:butterfly}{kepompong}: kebanyakan \omterm{prop:g3:sorting-living-things:groups}{serangga}, seperti belalang atau kupu-kupu.
\end{solution}

\begin{solution}{exo:g6:life-through-seasons:7}
Kupu-kupu: sebuah \omterm{ex:g2:animals-grow-up:butterfly}{kepompong} di bawah batu puncak tembok. Belalang:
\omterm{def:g2:animals-grow-up:born}{telur} di dalam tanah. Kepik: seekor dewasa yang \omterm{prop:g1:healthy-habits:needs}{tidur} berimpitan jauh
di dalam serasah. Siput: tersegel di dalam \omterm{def:g1:animals-around-us:covering}{cangkangnya} di tebing
tanahnya. \omterm{prop:g3:sorting-living-things:groups}{Burung} layang-layang: di Afrika, sebagai dirinya yang
\omterm{def:g3:animals-through-seasons:migration}{bermigrasi}.
\end{solution}

\begin{solution}{exo:g6:life-through-seasons:8}
Supaya setiap perbedaan di antara keempat tabelnya datang dari
musimnya dan bukan dari penyurveiannya --- aturan uji yang adil pada
\cref{met:g6:exploring-our-environment:survey}, yang diterapkan
melintasi waktu alih-alih melintasi lokasi.
\end{solution}

\begin{solution}{exo:g6:life-through-seasons:9}
Pertama, banyak ``yang absen'' sebenarnya hadir dalam wujud tersembunyi
--- \omterm{def:g2:animals-grow-up:born}{telur}, \omterm{ex:g2:animals-grow-up:butterfly}{kepompong}, dewasa yang \omterm{prop:g1:healthy-habits:needs}{tidur}, \omterm{ex:g4:plants-without-seeds:bulbtuber}{umbi lapis} --- yang
dilewatkan sebuah survei berjalan-jalan tanpa penggalian dan
pengangkatan. Kedua, sebagian tidak hilang melainkan berada di tempat
lain: para \omterm{def:g3:animals-through-seasons:migration}{pemigrasi} akan kembali. Daftar bulan Januari meremehkan
jumlah populasinya dan menghitung para pengembara sebagai lenyap.
\end{solution}

\begin{solution}{exo:g6:life-through-seasons:10}
\omterm{def:g1:plants-around-us:plant}{Tumbuhannya} berada dalam keadaan gudang bawah tanahnya --- hidup, dan
menunggu sebagai \omterm{ex:g4:plants-without-seeds:bulbtuber}{umbi lapis}. Pada bulan Maret umbinya akan
menumbuhkan \omterm{def:g1:plants-around-us:parts}{daun} dan \omterm{def:g3:flowers-fruits-seeds:parts}{bunga} lagi, lebih rapat kalau rumpunnya sudah
berlipat ganda.
\end{solution}

\begin{solution}{exo:g6:life-through-seasons:11}
Hilang ke mana, sebagai apa? Katak melewati musim dingin dalam keadaan
dorman --- terkubur di dasar lumpur kolamnya atau di tempat berlindung
yang lembap di darat --- melewatkan bulan-bulan yang dingin dalam
keadaan lesu sampai musim semi memanggilnya kembali ke air untuk
berbiak.
\end{solution}

\begin{solution}{exo:g6:life-through-seasons:12}
Masalah musim dingin bagi sehelai \omterm{def:g1:plants-around-us:parts}{daun} yang bekerja adalah menahan air:
tanah yang membeku memberi \omterm{def:g1:plants-around-us:parts}{akarnya} nyaris tidak ada air, sementara
angin dan matahari terus menariknya keluar --- masalah pengeringan
kutu kayu, sepanjang musim dingin. Lilin dan bangun jarum pada \omterm{def:g1:plants-around-us:parts}{daun}
yang \omterm{def:g6:life-through-seasons:trees}{selalu hijau} memangkas kehilangan itu sampai ke ukuran yang
sanggup ditanggung pohonnya. \omterm{def:g1:plants-around-us:parts}{Daun} yang lebar dan tipis tidak dapat
disegel sebaik itu; jawaban pohon \omterm{def:g6:life-through-seasons:trees}{meranggas} yang lebih murah adalah
menghapus \omterm{def:g1:plants-around-us:parts}{daun} tipisnya tiap musim gugur lalu menyimpan musim seminya
di dalam \omterm{def:g6:life-through-seasons:buds}{kuncup}.
\end{solution}

\begin{solution}{pb:g6:life-through-seasons:1}
\textbf{1.} \omterm{def:g6:life-through-seasons:buds}{Kuncup} pada setiap \omterm{def:g1:plants-around-us:tree}{rantingnya} --- gudang musim dingin yang
menyimpan tunas musim semi berikutnya, yang dibangun pada musim panas
sebelumnya.

\textbf{2.} Di antara bingkai bulan Maret dan April (dengan hari yang
memanjang dan pagi yang lebih lunak pada catatannya). Pohonnya
menjawab \omterm{def:g6:exploring-our-environment:conditions}{suhu} yang naik dan \omterm{def:g6:exploring-our-environment:conditions}{cahaya} yang memanjang.

\textbf{3.} \omterm{prop:g3:sorting-living-things:groups}{Serangga} penyerbuk --- terutama lebah --- harus
mengunjungi \omterm{def:g3:flowers-fruits-seeds:parts}{bunganya}, sehingga \omterm{def:g3:flowers-fruits-seeds:parts}{serbuk sari} mencapai \omterm{def:g3:flowers-fruits-seeds:parts}{putiknya} dan
\omterm{def:g5:flower-to-fruit:steps}{pembuahan} menyusul. \omterm{prop:g3:flowers-fruits-seeds:transform}{Buah} berangan adalah \omterm{def:g2:from-seed-to-plant:seed}{biji} yang terbentuk dari
sebuah \omterm{prop:g5:flower-to-fruit:roles}{bakal biji} yang terbuahi, di dalam selongsong \omterm{prop:g3:flowers-fruits-seeds:transform}{buahnya} yang
berduri.

\textbf{4.} \omterm{def:g6:life-through-seasons:trees}{Meranggas}. Musim semi berikutnya menumpang bingkai musim
dinginnya sebagai \omterm{def:g6:life-through-seasons:buds}{kuncup} pada \omterm{def:g1:plants-around-us:tree}{ranting} yang telanjang.

\textbf{5.} Di bawah halaman rumputnya, sebagai \omterm{ex:g4:plants-without-seeds:bulbtuber}{umbi lapis} --- gudang
bawah tanah milik \omterm{prop:g6:life-through-seasons:herbs}{tumbuhan tahunan}.

\textbf{6.} Sebagai \omterm{def:g2:from-seed-to-plant:seed}{biji} di dalam tanah: siasat \omterm{prop:g6:life-through-seasons:herbs}{tumbuhan semusim} ---
\omterm{def:g1:plants-around-us:plant}{tumbuhannya} sendiri mati pada musim gugur.

\textbf{7.} \omterm{prop:g3:sorting-living-things:groups}{Burung} layang-layang: di Afrika, sebagai dirinya yang
\omterm{def:g3:animals-through-seasons:migration}{bermigrasi}, mengikuti \omterm{prop:g3:sorting-living-things:groups}{serangga} terbang. Lebahnya: di dalam sarang,
sebagai gerombolan musim dingin, hangat mengelilingi ratunya.

\textbf{8.} \omterm{def:g3:animals-through-seasons:active}{Tetap aktif}. \omterm{prop:g3:sorting-living-things:groups}{Burung} murainya harus menemukan makanan
sepanjang bulan-bulan yang paling kurus --- persis masalah yang
dihindari \omterm{prop:g3:sorting-living-things:groups}{burung} layang-layang dengan pergi.

\textbf{9.} Misalnya: \omterm{def:g6:life-through-seasons:buds}{kuncupnya} merekah (Maret--April) ketika harinya
mencapai $11$--$13$ jam dan paginya menghangat melampaui
\qty{8}{\degreeCelsius}; \omterm{def:g1:plants-around-us:parts}{daunnya} menguning keemasan lalu gugur
(Oktober--November) ketika harinya memendek dan paginya mendingin.
Peristiwa pohonnya berpasangan dengan perubahan \omterm{def:g6:exploring-our-environment:conditions}{kondisinya}, bingkai
demi bingkai.

\textbf{10.} Petaknya tidak mati: ia menarik diri ke gudangnya ---
hidup sebagai \omterm{ex:g4:plants-without-seeds:bulbtuber}{umbi lapis} di bawah tempat yang sama. ``Hilang'' selalu
pantas menerima: hilang ke \emph{mana}, sebagai \emph{apa}?

\textbf{11.} Dahan telanjang yang membawa \omterm{def:g6:life-through-seasons:buds}{kuncup}; seekor \omterm{prop:g3:sorting-living-things:groups}{burung} murai
yang mungkin ada di suatu tempat; di bawahnya, serasah dan tanah \omterm{def:g3:flowers-fruits-seeds:parts}{bunga}
narsis yang kosong tetapi menyembunyikan \omterm{ex:g4:plants-without-seeds:bulbtuber}{umbi lapis}; \omterm{def:g2:from-seed-to-plant:seed}{biji} \omterm{def:g3:flowers-fruits-seeds:parts}{bunga} popi
di dalam tanah jalan setapaknya. Taruhan yang pasti adalah \omterm{def:g6:life-through-seasons:buds}{kuncup} dan
ketelanjangannya --- keduanya mengikuti dari permesinan \omterm{def:g6:life-through-seasons:trees}{meranggas}
pohonnya dan bukti setahun; \omterm{prop:g3:sorting-living-things:groups}{burung} murainya mungkin ada, tidak
dijanjikan.

\textbf{12.} Ini menunjukkan bahwa kehangatanlah \omterm{def:g6:exploring-our-environment:conditions}{kondisi} yang hilang
--- \omterm{def:g6:exploring-our-environment:conditions}{cahaya} di dalam ruangan pada bulan Februari tidak lebih panjang
daripada di luar, tetapi air dan ruangannya hangat, dan itu sudah
cukup. Uji yang adil: dua \omterm{def:g1:plants-around-us:tree}{ranting} dari pohon yang sama, keduanya di
dalam air di jendela yang sama, satu di ruangan yang hangat, satu di
serambi yang dingin --- satu perbedaan, yaitu kehangatan; kalau hanya
\omterm{def:g1:plants-around-us:tree}{ranting} yang hangat yang berdaun, putusannya sudah dipastikan.
\end{solution}
